% =====================================================================
% exam-cram — STUDY NOTES PDF (multi-page, readable revision notes)
% Compile:  tectonic study-notes.tex
%
% This is the "复习材料" deliverable: a properly typeset, multi-page
% LaTeX document generated FROM the ingested lecture/homework material.
% Unlike the 2-page cheat sheet, this one is READABLE: larger type,
% chapters per topic, display formulas, worked examples, and callouts.
%
% AUTHORING RULES (the model that fills this in must follow):
%  * One \section per topic (or major exam theme), \subsection for parts.
%  * Use \defn{...} for definitions, \formula{...} for display formulas,
%    \example{...} for worked examples, \warn{...} for common mistakes.
%  * Everything must be grounded in the ingested sources; add source tags
%    like (L1 p.4), (Slide 7), (HW2 Q3) at the end of items.
%  * Keep the preamble unchanged; fill in \title/\date and the body.
%  * Aim for one to several pages per topic; total length is natural
%    (5–20 pages typical). No page-fitting is applied.
% =====================================================================
\documentclass[11pt,a4paper]{article}

\usepackage[margin=2.2cm]{geometry}
\usepackage{amsmath,amssymb,mathtools}
\usepackage{enumitem}
\usepackage{booktabs}
\usepackage{array}
\usepackage[table]{xcolor}
\usepackage{fancyhdr}
\usepackage{titlesec}
\usepackage{fontspec}
\usepackage{xeCJK}
\usepackage{parskip}

\setCJKmainfont[BoldFont={Microsoft YaHei UI Bold}]{Microsoft YaHei}
\setCJKsansfont{Microsoft YaHei UI}

% ---- colours ----
\definecolor{navy}{HTML}{17324D}
\definecolor{acc}{HTML}{1F5C99}
\definecolor{ok}{HTML}{1F6B45}
\definecolor{bad}{HTML}{B3261E}
\definecolor{muted}{HTML}{5A6270}
\definecolor{defnbg}{HTML}{E8F0F8}
\definecolor{warnbg}{HTML}{FDECEA}
\definecolor{exbg}{HTML}{F2EAF8}
\definecolor{tipbg}{HTML}{E8F5EE}

% ---- heading styles ----
\titleformat{\section}{\Large\bfseries\color{navy}}{\thesection.}{0.6em}{}
\titleformat{\subsection}{\large\bfseries\color{acc}}{\thesubsection.}{0.5em}{}
\titleformat{\subsubsection}{\normalsize\bfseries\color{ink!85}}{\thesubsubsection.}{0.5em}{}
\titlespacing*{\section}{0pt}{14pt plus 2pt}{6pt}
\titlespacing*{\subsection}{0pt}{10pt plus 2pt}{4pt}

% ---- header/footer ----
\pagestyle{fancy}
\fancyhf{}
\fancyhead[L]{\small\color{muted}\leftmark}
\fancyhead[R]{\small\color{muted}\rightmark}
\fancyfoot[C]{\small\thepage}
\renewcommand{\headrulewidth}{0.3pt}

% ---- reusable blocks ----
% definition
\newcommand{\defn}[1]{\par\vspace{3pt}\noindent
  {\colorbox{defnbg}{\textcolor{acc}{\sffamily\bfseries 定义 / Definition}}\hspace{3pt}}%
  \textit{#1}\par}
% warning / common mistake
\newcommand{\warn}[1]{\par\vspace{3pt}\noindent
  {\colorbox{warnbg}{\textcolor{bad}{\sffamily\bfseries 易错 / Common error}}\hspace{3pt}}%
  #1\par}
% tip
\newcommand{\tip}[1]{\par\vspace{3pt}\noindent
  {\colorbox{tipbg}{\textcolor{ok}{\sffamily\bfseries 技巧 / Tip}}\hspace{3pt}}%
  #1\par}
% example
\newenvironment{example}{%
  \par\vspace{4pt}\noindent
  \begin{minipage}{\dimexpr\linewidth-2\fboxsep-0.8em}%
  \colorbox{exbg}{\textcolor{acc}{\sffamily\bfseries 例题 / Example}}\hspace{3pt}%
  \ignorespaces
}{\end{minipage}\par\vspace{2pt}}

% ---- document ----
\begin{document}

% =====================================================================
% TITLE
% =====================================================================
\begin{center}
  {\LARGE\bfseries\color{navy} 课程名 —— 复习材料}\\[2pt]
  {\large\color{muted} 考试：考试名 · 日期：YYYY-MM-DD · 时长/形式}\\[6pt]
  {\small\color{muted} 依据讲义与作业整理 · 来源标注 (L1 p.4) = 讲义1第4页, (Slide 7), (HW2 Q3) 等}
\end{center}
\vspace{4pt}
{\color{navy}\hrule height 1pt}\vspace{6pt}

\tableofcontents
\clearpage

% =====================================================================
% TOPIC 1
% =====================================================================
\section{主题一：名称}

% 总体一句话
\begin{quote}
  <本主题一句话概括：它解决什么问题、考试里怎么出现>
\end{quote}

\subsection{核心概念}
\defn{<定义1>（(L1 p.4)）}
\defn{<定义2>（(Slide 3)）}

<解释性段落，用自己的话把讲义里的逻辑讲清楚。可以分多段，每段 2–4 行。>

\subsection{公式与条件}
\begin{align}
  <公式1> \quad \text{条件：} <条件> \label{eq:1}
\end{align}

\begin{align}
  <公式2> \quad \text{适用：} <何时用> \label{eq:2}
\end{align}

\tip{<快速记忆/检查技巧>}

\subsection{典型例题（来自讲义/作业）}
\begin{example}
\textbf{例 1（(HW2 Q3)）：}<题目描述>

\textbf{解：}<分步解答，最后给出答案形式>
\end{example}

\subsection{常见错误}
\warn{<错误1 及正确做法>}
\warn{<错误2 及正确做法>}

% =====================================================================
% TOPIC 2 — 依此类推，每个主题一节
% =====================================================================
\section{主题二：名称}
...

% =====================================================================
% FINAL: 公式总表 / 复习清单
% =====================================================================
\section*{附：公式总表}
\begin{center}
\begin{tabular}{@{}p{0.32\linewidth}p{0.32\linewidth}p{0.32\linewidth}@{}}
\toprule
\textbf{公式} & \textbf{适用条件} & \textbf{常见错误} \\
\midrule
$<公式>$ & $<条件>$ & $<错误>$ \\
... & ... & ... \\
\bottomrule
\end{tabular}
\end{center}

\end{document}